\Lampiran{Contoh Tabel Lampiran}

\noindent
Tuliskan isi lampiran pertama di sini, misalnya tabel data lengkap.

\begin{center}
  \renewcommand{\arraystretch}{1.2}
  \begin{tabular}{|c|c|c|}
    \hline
    \textbf{No} & \textbf{Data} & \textbf{Keterangan} \\ \hline
    1 & \dots & \dots \\ \hline
    2 & \dots & \dots \\ \hline
  \end{tabular}
\end{center}

\Lampiran{Contoh Kode Program}

\noindent
Kode program dapat dituliskan menggunakan lingkungan \texttt{lstlisting}.

\begin{lstlisting}[language=Python, style=mystyle]
def iterasi(T, x0, tol=1e-6):
    # Iterasi titik tetap sederhana
    x = x0
    while True:
        x_baru = T(x)
        if abs(x_baru - x) < tol:
            return x_baru
        x = x_baru
\end{lstlisting}
